\section{System: Hard Verification, Abstention, and the Movable Projection Boundary}\label{sec:chapter-5-system-hard-verification-abste}

Chapter 3 built the knowledge graph, and Chapter 4 built the benchmark that measures question answering over it. This chapter describes CICADA, the system that answers those questions over the same graph. Its central idea is a strict division of labour: a language model proposes, deterministic code checks and executes, and the system abstains where neither can support an answer.

\section{Design goal and the three-layer principle}\label{sec:design-goal-and-the-three-layer-principl}

CICADA's design goal follows from the constraint of Chapter 1: an answer is only ever produced by deterministic execution against the knowledge graph, and where the system cannot support an answer it must abstain. Between question and executed answer lies a chain of semantic judgements: what kind of answer is wanted, which literals constrain it, which organisation plays which role, and whether the concept asked about exists in the graph at all. A language model is indispensable for making these judgements and untrustworthy for certifying them. The architecture therefore asks one question of every property on that chain: \emph{can it be projected onto structure or execution?} If it can, it is pinned by a deterministic check, and the model is never trusted for it. If it cannot, it falls into a blind region that no mechanical check can see into. The response to that region is twofold. Calibrated abstention catches what projection misses, and learning --- trained only on signal from the verifiable region --- gradually covers part of what remains.

This yields the three layers that structure the chapter. The \textbf{hard-verification layer} (\S5.3--5.4) pins every projectable property: compilation, grounding, execution, and structural cross-checks between question, briefing, and plan. The \textbf{abstention layer} absorbs the unprojectable residue; its completeness is tested by the benchmark's three abstention-trap families (ambiguous, schema-unsupported, and empty-result questions; Chapter 4). The \textbf{learning layer} (Chapter 6) consumes only what the first layer certifies. The orchestrator records this division of authority in every trace it emits, as a literal string in the workflow metadata: \emph{"LLM proposes; KG executes; verifier judges; reflector refines"} (\texttt{pipeline.py}, \texttt{\_workflow\_trace}). Two further claims give the chapter its shape. First, the boundary between the projectable and the unprojectable is \emph{movable}: \S5.6 documents one move, the v2 revision, including the measured cost of moving it too far. Second, the reflector that repairs failed plans sits deliberately \emph{outside} the verification chain. It consumes verifier signal but never joins the chain, and every repair it proposes is re-checked by that chain (\S5.5).

\begin{figure}[t]
\centering
\includegraphics[width=\linewidth]{architecture}
\caption{The CICADA architecture. A question passes through two local LoRA planners (Step~1 briefing, Step~2 typed graph plan) into the deterministic core, the sole answer authority. Circled numbers mark the four-check chain; check~3 (briefing$\leftrightarrow$question) is the semantic seam where hard negatives leak. The reflector consumes diagnostic signals and its repairs re-enter all four checks. Verified traces alone supervise the two adapters (learning layer).}
\label{fig:architecture}
\end{figure}

\section{Two-step planning: the language model as translator}\label{sec:two-step-planning-the-language-model-as}

The language model appears in exactly one role: translating the question into a machine-checkable intermediate representation. It never sees the graph's contents, never composes an answer, and never adjudicates its own output. Everything it emits is a \emph{claim about the question}, expressed in a vocabulary small enough for deterministic code to audit.

\textbf{Step 1: the intent-program briefing.} The first call (\texttt{question\_intent\_program\_messages}, \texttt{typed\_planning.py}) asks a reading model for a structured \emph{briefing}. The briefing contains an \texttt{answer\_signature} (operation, value type, natural-language answer-field description), a \texttt{program} of at most eight typed steps whose operations come from a closed enumeration (\texttt{filter\_records}, \texttt{bind\_filter}, \texttt{distinct}, \texttt{count}, \texttt{sum}, \texttt{select}, \texttt{exists}, \texttt{compare}, \texttt{argmin}/\texttt{argmax}, \texttt{top\_k}, \texttt{abstain}), an \texttt{answer\_step}, and an \texttt{unsupported\_or\_ambiguous} list. The prompt's hard rules are contractual: root steps take no inputs (the model may not invent a universe table); literals must be copied verbatim, never invented; and an unsupported or ambiguous question must be expressed as \texttt{operation=abstain} with an empty program. Filters describe fields in natural language, such as "published year" or "buyer organisation". The mapping to schema columns is left to the deterministic grounding layer (\S5.3). A well-formed intent program takes a \textbf{fast path}: it compiles deterministically into a graph plan without a second model call, subject only to a narrowly targeted veto (\texttt{\_fastpath\_veto}, \S5.6). Malformed programs are not reinterpreted. A contract verifier (\texttt{\_verify\_intent\_program}) checks every step's operation against its declared return type, every input against earlier steps, and the answer step against the signature; violations become structured failure codes rather than guesses.

\textbf{Step 2: the typed graph plan.} When the fast path is unavailable, a second call (\texttt{typed\_plan\_messages}, lean variant) turns the question plus briefing into a \texttt{graph\_plan}. This is a flat array of variables --- each with a kind (\texttt{entity}, \texttt{entity\_set}, \texttt{record\_set}, \texttt{scalar}), a role, and filters over an enumerated slot vocabulary --- plus relations with typed bind fields and an enum-closed \texttt{return} object. Variable identifiers encode structure (\texttt{a1} feeds the answer, \texttt{b1a1} feeds \texttt{a1}), so the dependency DAG is readable from a flat array even when the model omits explicit edges. Prompt and schema variants are resolved per model from measured ablations (\texttt{resolve\_planner\_variants}); the local 8B students take the lean prompt with the omit-unused schema.

\textbf{Guided decoding as format authority.} Both steps are shape-enforced by strict JSON schemas (\texttt{intent\_program\_schema}, \texttt{graph\_plan\_schema}) passed to the provider or, for local models, to vLLM's guided decoding. Enumerations do real work here: slot names, operations, and roles are enum-constrained, which eliminates placeholder echo and invented vocabulary at the decoding level. Format is therefore externally guaranteed for every system evaluated, including the untrained zero-shot rung. Differences along the training ladder measure planning quality, not format compliance; \S5.7 gives the diagnostic that makes this point empirically.

\textbf{Structural resampling.} The evaluation configuration runs \texttt{plan\_samples=2}. If the first sample fails compilation or the consistency check --- a \emph{detectable} defect with a structured reason --- the planner requests one fresh, deliberately different decomposition before falling back to abstention or repair. The boundary is principled: an executable plan that is semantically wrong is never resampled, because at inference time no oracle exists to detect it.

\section{The deterministic core: compile, ground, execute, attest}\label{sec:the-deterministic-core-compile-ground-ex}

Everything downstream of the planner is deterministic code, and it is deliberately the largest part of the system.

\textbf{Compilation and structural repair.} \texttt{compile\_graph\_plan} (\texttt{graph\_planning.py}) converts the model's graph JSON into an executable plan. It derives dependencies from the identifier convention and treats relations as dataflow edges, because a relation without a matching dependency would otherwise silently never bind; explicit edges win over naming-derived ones when they conflict. Before compilation, \texttt{normalise\_graph\_plan} applies a catalogue of \emph{T-transforms}: deterministic rewrites of mechanically recognisable planner idioms into the structure they evidently mean, never guessing semantics. Among them: a filter whose value names another variable becomes a dependency edge (T1). An empty-valued filter is schema padding, never executable, and is dropped (T2). A compare side that is a scalar variable holding one literal amount or date is folded into the return as a threshold, rather than executed as a spurious record filter (T11/T12). An echoed bridge anchor (the organisation both bound and literally re-filtered) is stripped so the bind carries the meaning (T9). A singular interrogative ("which buyer\ldots{}") answered with \texttt{distinct\_set} is rewritten to \texttt{select}, so run-time uniqueness adjudicates instead of a 30-item list hiding an ambiguity (T10).

Compilation itself enforces DAG validity, with a second-chance compile that drops naming-derived edges when they alone create a cycle. It rejects an \emph{unconstrained bind source}: an entity set with no filters and no dependencies is the whole dimension, and binding on it would silently answer a different question. It also grounds \emph{dangling} variables, so an unsupported concept smuggled into a never-executed variable still fails the plan with its reason intact. Forms the single-query executor cannot express are rewritten. A two-organisation comparison becomes a \texttt{DecompositionPlan} of two independent sub-counts combined by \texttt{compare\_gt}; \texttt{rank\_top\_k} maps onto the grouped \texttt{top\_k}; and \texttt{find\_extreme} maps onto \texttt{argmax}/\texttt{argmin} over the exhaustive match set, answering with the extremum \emph{record} (its \texttt{contract\_node\_id}) rather than a bare value.

\textbf{Schema grounding.} Natural-language field descriptions from Step 1 map to canonical slots via a two-stage grounder (\texttt{schema\_grounding.py}). Candidates are generated over an alias table plus a character n-gram embedding; validation then applies a type gate, a confidence threshold of 0.72, and a top-2 margin of 0.08. An ambiguous match is a structured failure, not a coin flip. Two asymmetries matter. First, a direction cue outranks the role noun: "buyer has awarded a contract to X" contains \emph{buyer}, but the "awarded \ldots{} to" phrasing names X as the supplier, and the cue wins. Second, a field text containing a token from the unsupported-concept lexicon (\texttt{\_UNSUPPORTED\_FIELD\_TOKENS}: invoice, payment, delivery, performance, reliable, bidder, fairness) is denied the permissive substring shortcut. "Invoice or payment date" can therefore no longer ground to the date slot and be answered from the award date --- the v2.2 rule whose calibration history \S5.6 recounts.

\textbf{Entity resolution.} Organisation surfaces resolve under a deliberately conservative policy (\texttt{resolve\_confident\_org}). The top candidate is accepted only if it is an exact record match, or if it scores at least 0.85 with a margin of at least 0.10 over the runner-up. Near-ties are inspected for \emph{decoration variance} --- "The X", "X (BEIS)", ampersand and case variants of one organisation. Because the organisation's rows are split across those surfaces, such a mention is resolved as an \texttt{in} filter over every variant rather than a silent pick of one. Genuinely distinct near-ties fail as \texttt{ambiguous\_entity\_candidates}. No caller anywhere takes \texttt{hits[0]} under a weak match.

\textbf{Spec grounding.} Immediately before execution, \texttt{ground\_spec} (\texttt{grounding.py}) rebuilds every field, operation, and guard of the proposed query. The executable operation set is closed (\texttt{select\_unique}, \texttt{count}, \texttt{sum}, \texttt{exists}, \texttt{argmax}, \texttt{argmin}, \texttt{distinct\_set}, \texttt{set}, \texttt{top\_k}, \texttt{predicate}); requests outside it fail with a stated reason (averages unsupported; \texttt{compare} and \texttt{in\_subquery} routed to the decomposition planner). Field aliases snap to canonical columns. Mechanically impossible value shapes are rejected \emph{because they would otherwise succeed}: an \texttt{in} filter over free text, or a bare year against a timestamp column, executes "successfully" and returns a verified-looking zero. Every money aggregation --- a sum, a predicate over a sum, a \texttt{top\_k} ranked by summed value --- has a \texttt{value\_is\_additive} guard added if the planner omitted it, so a total never silently includes non-additive framework-ceiling values. Every reduction is forced onto exhaustive retrieval.

\textbf{Execution.} The executor (\texttt{executor.py}) runs the grounded spec over the full knowledge graph, never a retrieved sample, after preflight checks on schema fields, constraint contradictions, and exhaustiveness. Match sets are deduplicated on \texttt{contract\_node\_id}. \texttt{count} and \texttt{exists} answer over the possibly empty match set, so 0 and False are legitimate verified answers. \texttt{select\_unique} returns \texttt{multiple\_answers} when the match set disagrees, which becomes a clarification rather than an enumeration. \texttt{top\_k} grouping excludes empty-string keys, so unnamed populations cannot occupy a rank. Monetary arithmetic is exact \texttt{Decimal} arithmetic with non-finite values barred. Evidence is attested exactly: the bundle carries the exact match count, with row identifiers capped at 200 and the cap flagged, so an audit never confuses a sample with a population.

\textbf{Evidence and answer cards.} A passed execution is wrapped in a KG-first evidence verdict (\texttt{evidence.py}) binding the claim to evidence identifiers. Post-flight disclosures never change the answer but make it self-explaining: a count discloses how many matched contracts lack a recorded supplier, and a factoid confirmed across \emph{n} agreeing contracts reveals hidden multiplicity. A failed or unsupported execution produces an abstaining answer card carrying the structured reason. Abstention on emptiness is \emph{provenance-gated}: an empty result whose every literal traces to the question is treated as the graph's answer, a rule enforced at the repair gate (\S5.5).

\section{The four-check chain and the seam}\label{sec:the-four-check-chain-and-the-seam}

The machinery of \S5.3 can be summarised as four checks on the correspondence chain from answer back to question. Their character differs, and stating the difference precisely is the centre of this chapter's argument.

\textbf{Check 1 --- executor $\leftrightarrow$ plan graph: hard.} The answer is produced \emph{by} compiling and executing the plan. The correspondence between what the plan says and what the executor does holds by construction, deterministically, on every question.

\textbf{Check 2 --- plan graph $\leftrightarrow$ briefing: hard.} The briefing's answer signature, declared template, and program structure are compared against the compiled plan in code. The intent-program contract verifier rejects an answer step whose operation or return type contradicts the signature, and the consistency check rejects an operation outside the declared question type. This is structural comparison, exhaustive over what it covers.

\textbf{Check 3 --- briefing $\leftrightarrow$ question: the seam.} Here code can check only \emph{projections} of meaning, never meaning itself. \texttt{plan\_consistency\_check} audits every projection the question text affords. Literal filter surfaces must occur in the question (\texttt{surface\_not\_in\_question}). Every number in a constraint or threshold must be question-supported, with deterministic decompositions for scaled amounts ("\pounds{}2.5m" supports 2500000) and written dates (\texttt{invented\_number}, \texttt{invented\_threshold}). Buyer/supplier roles are checked against direction-cue windows around the mention, with bridge anchors exempted (on "suppliers who worked with X" the plan \emph{correctly} constrains X as buyer) and interrogative answer-role phrases masked (\texttt{role\_flipped}). Explicit atoms (years, CPV codes, categories, quoted titles) must survive into the compiled spec (\texttt{missing\_explicit\_constraint}). Strong bridge phrasings on a genuinely flat plan flag a collapsed intermediate set (\texttt{bridge\_cue\_requires\_bridge\_join}, fired only when the graph carries no dependency structure). Comparison direction words must support the plan's comparator. A cue lexicon of concepts the graph does not carry forces the unanswerable reading. A record-count question whose return edge counts an \texttt{entity\_set} variable is rejected as the count-intermediate-set misread (\S5.6). Each of these is a hard check \emph{of a projection}. None of them, and no possible extension of them, checks whether the plan means what the question means. Check 3 is therefore the entrance to the blind region, and exactly here \textbf{hard negatives} arise: plans that pass every projection, ground, execute, and verify cleanly, and are nevertheless wrong. When the training harvest later compares verifier-passing outputs against the offline oracle, the verifier-passed-but-oracle-wrong residue is exactly the population that leaked through this seam. Those specimens are collected as \texttt{hard\_negatives}; their use and misuse as preference-learning material is the subject of Chapter 7.

\textbf{Check 4 --- answer $\leftrightarrow$ oracle: hard, but training-time only.} During data harvest, and only there, a fourth deterministic check compares the executed answer against an independently computed oracle. The oracle filters candidate training examples; it never authors them. It does not exist in the inference pipeline, and it must not: consulting it at evaluation time would be cheating, and every number reported in this dissertation is produced without it.

The framing this dissertation commits to has a negative half that matters as much as the positive half. Everything hard-verifiable is hard-verified; what cannot be hard-verified is compressed to a minimum; abstention and learning catch the remainder. The four checks squeeze the blind region layer by layer, but they do not eliminate it.

\section{Gated repair: the reflector as a consumer of verifier signal}\label{sec:gated-repair-the-reflector-as-a-consumer}

When a check fails, the failure is not the end of the trace. A second language-model component --- the \emph{reflector} --- receives it and proposes a repaired plan. What the reflector is not is essential: it is \textbf{not a fifth check}, and this dissertation never describes it as a verifier. The four deterministic checks \emph{produce} diagnostic signal. The reflector, an uncertain model, only \emph{consumes} that signal to propose a repair, and the proposal re-enters the same four checks. The reflector never certifies its own fix: a repair is accepted only when the full ground $\rightarrow$ execute $\rightarrow$ verify path passes it. The design deliberately refuses the alternative, an uncertain model verifying an uncertain model, because that pattern invites an endless regress.

The signal the reflector consumes is rich and strictly answer-free. \texttt{\_feedback\_from\_trace} assembles the failure stage (grounding, schema, executor, sanity, verifier, or plan-compile), the structured reason, the failed plan with its raw graph surfaced even on compile failures, the failed variable, and grounding issues with top candidates. A scrubber strips anything oracle-shaped before any model sees it. A deterministic reason-to-prescription table (\texttt{\_REPAIR\_GUIDELINES}) prepends targeted guidance: a cycle diagnosis names which edge direction to keep, and \texttt{no\_results\_with\_invented\_literals} says to remove exactly the literal the question never stated. The reflector answers under its own strict schema: a diagnosis, one action from a closed set of six (\texttt{fix\_question\_type}, \texttt{repair\_constraints}, \texttt{swap\_buyer\_supplier}, \texttt{change\_operator}, \texttt{change\_operation}, \texttt{abstain}), the repaired plan, and explicit changed/unchanged field lists.

Repair is \emph{gated} before it is attempted. The eligibility rule (\texttt{\_replan\_eligibility}) separates a repairable defect from a legitimate abstention. A plan that never executed, an executor error, or a \texttt{no\_results} caused by a planner-invented literal may be repaired. But a **clean \texttt{no\_results} --- grounded, executed, every user-visible literal traceable to the question --- is the knowledge graph's answer**, and repairing it is answer-shopping. The provenance test (\texttt{\_invented\_literals}) excludes bound values, invisible guards, and long derived in-lists, so only genuinely invented user-facing literals open the gate. The motivating measurement comes from the $\leq$50-question development smoke slice and is cited as directional only. A naive "no answer $\rightarrow$ replan" loop \emph{reduced} accuracy from 84\% to 66\% and raised hallucinated answers from 1 to 6, repairing six correct abstentions into confident wrong answers. The gate exists to make repair monotone.

The loop is bounded and fully recorded as an \textbf{attempt protocol}. Attempt 0 is the initial plan; each subsequent attempt consumes the \emph{latest} failure's feedback; and every attempt's plan, status, and feedback lands in \texttt{metadata["feedback\_replan"]["attempts"]}. This record is the substrate of the learning layer. A harvest run computes Repair@k from it and mines preference pairs directly (chosen = the first verified attempt, rejected = the nearest preceding failure), which is where the DPO training data of Chapter 6 originates. Two repair budgets exist as deliberately distinct configurations, never to be conflated. The \textbf{evaluation protocol} runs \texttt{max\_feedback\_replans=1} (\texttt{run\_compare.py}), so every ladder rung is scored under the same one-repair budget. The \textbf{harvest protocol} (\texttt{run\_teacher.py}) runs \texttt{--max-repairs 2}, spending a larger budget to maximise verified yield. Evaluation numbers are never numbers "under the harvest config", and the two metrics they feed --- system accuracy versus data-engine yield --- are kept apart throughout.

Harvest routing completes the interface to the learning layer. Verifier-and-oracle-passed traces become \texttt{verified\_sft}; verifier-passed-but-oracle-wrong traces become \texttt{hard\_negatives} (the seam's leakage made visible); correct abstentions become \texttt{abstain\_sft}; successful repairs become \texttt{repair\_sft}; and the attempt protocol yields preference pairs. Chapter 6 takes over from there.

\section{The movable projection boundary: the v2 revision as a worked example}\label{sec:the-movable-projection-boundary-the-v2-r}

The blind region is not a fixed set. The seam is where projection \emph{currently} stops, and engineering effort can move it: a property that today only the oracle can see may tomorrow be a structural check. The v2 revision shows this in practice, in both directions.

The revision began from an error autopsy of the v1 development matrix. compare\_v4 is the development set; every fix below was derived from its v1 errors and confirmed only later on the untouched final test set. The autopsy found 21 abstentions on answerable questions, 12 answers on unsupported questions, and 22 wrong values, root-caused to four precise defects. Each fix is a projection move: a semantic error class re-expressed as a structural property.

\begin{enumerate}
\item \textbf{The count-target check.} The dominant verifier-invisible bridge failure was the \emph{count-intermediate-set} misread. On "how many contract notices\ldots{}" questions, the plan counted the intermediate entity set instead of the target records; in the motivating trace, 30 CPV codes where the oracle counted 1,439 notices. Such a plan is semantically wrong, structurally flawless, and invisible to execution. \texttt{\_count\_target\_is\_entity\_set} makes the defect structural: a record-count question whose return edge counts an \texttt{entity\_set} variable is rejected before execution. The trigger is deliberately precise --- "how many CPV codes/suppliers" does not match the record-count pattern.
\item \textbf{The empty-value type-gate exemption.} Categorical lookups ("what is the procurement category of\ldots{}") ground a field text with an \emph{empty} value. v1's type gate rejected them for failing the goods/services/works enumeration, silently unmapping every categorical lookup's return field. The v2 gate exempts the empty value, because a lookup context has nothing to validate, while still rejecting any non-enum \emph{stated} value.
\item \textbf{The fast-path targeted veto.} The intent-program fast path (\S5.2) bypassed Step 2 entirely, and with it the two consistency checks that catch failures execution cannot: the unsupported-cue check and the count-target check. \texttt{\_fastpath\_veto} applies exactly those two checks to compiled fast-path graphs --- and, deliberately, \emph{only} those two.
\item \textbf{The unsupported-field-token rule.} Alias grounding's permissive substring rule let "invoice or payment date" ground to the date slot at high confidence, so an unsupported-concept question was answered from the award date. The v2.2 rule (\S5.3) blocks the substring shortcut exactly when the field text names a concept from the unsupported lexicon.
\end{enumerate}

\textbf{The iteration was not clean, and the record is kept honest.} The first attempt, v2.0, ran the \emph{full} consistency battery on fast-path graphs. Those checks are calibrated for raw planner payloads, and on compiled graphs they false-positive: the teacher's development leg regressed several points, with answerable questions flipping to null on cue-free wording. v2.1 narrowed the gate to the two-check veto above. The second over-reach was subtler. v2.1 replaced the substring alias rule with token-coverage scoring, which passed the entire synthetic unit battery --- and then rejected \textbf{88\% of 169 real field texts} harvested from actual teacher-trace briefings. Real field texts are sentences ("Was the award signed after 1 May 2025?") that v1's substring rule had grounded without difficulty. The teacher's factoid bucket collapsed, while the students, whose trained plans carry canonical slots, were unaffected. v2.2 reverted to v1's permissiveness \emph{except} for the unsupported lexicon, and was verified decision-identical to v1 on all 169 real field texts while still rejecting the unsupported-concept phrasings. Two lessons are recorded. First, \textbf{over-projection wounds the system exactly as under-projection does}. A check that fires on legitimate traffic converts answers into abstentions, just as a missing check converts abstentions into hallucinations. Moving the boundary therefore demands regression against \emph{real pipeline intermediates}, not synthetic unit cases. Second, single-run teacher deltas are treacherous. The subsequent three-replicate noise floor (v2.2 teacher 72.6\% $\pm$ 1.1, replicates 71.9/71.9/73.9, against a v1 single run of 70.0\%) showed that part of the v2.0/v2.1 regression chase was provider noise. Single-run teacher differences below about three points are therefore treated as unclaimable. The v2.0 and v2.1 intermediate matrices are archived as non-citable; only v1 and v2.2 are citable pipeline versions.

\textbf{The measured effect of the boundary move.} On the development set (n=260), v1 $\rightarrow$ v2.2 lifts every rung of the ladder by roughly five points: zero-shot 61.5\% $\rightarrow$ 70.4\%, SFT 76.1\% $\rightarrow$ 81.2\%, RSFT 76.1\% $\rightarrow$ 81.5\%, DPO 78.8\% $\rightarrow$ 83.5\%. The reading discipline is strict. This is a \emph{scaffolding} effect --- the same adapters, a raised floor for every system including the untrained zero-shot rung and the teacher --- and it must never be reported as a training gain. Training gains are measured on top of the raised floor, and Chapter 6 owns that decomposition. What the five points measure is the value of one boundary move: formerly semantic error classes are now caught structurally, before execution, for every system that runs the pipeline. What could not be carried across the line, plan-versus-question meaning in the general case, remains with abstention and with learning.

\begin{figure}[t]
\centering
\includegraphics[width=0.82\linewidth]{F3_pipeline_ablation}
\caption{Deterministic-scaffolding ablation: the same four adapters evaluated under pipeline v1 and v2.2. Every rung gains; the adapters are unchanged.}
\label{fig:ablation}
\end{figure}

\section{Serving: one GPU, guided decoding, fully local}\label{sec:serving-one-gpu-guided-decoding-fully-lo}

The deployment target motivates the final design decisions. Each student stage serves from a single local GPU via vLLM's OpenAI-compatible server (\texttt{serve\_student.sh}). The base model loads once, and fine-tuned stages are hot-plugged as named LoRA adapters (\texttt{--enable-lora --max-lora-rank 64}) with no weight merging. The zero-shot rung and every trained rung therefore share identical base weights and serving configuration (Qwen3 in non-thinking mode, matching its training distribution). Format enforcement is uniform at request time: clients send \texttt{response\_format=json\_schema}, which vLLM maps to guided decoding, so \emph{every} ladder rung --- including the untrained base --- is format-locked, and rung-to-rung deltas measure planning skill rather than JSON discipline. In the fully-local topology, the evaluation driver points Step 1 at a second local vLLM endpoint (\texttt{--step1-base-url}), and no request of any kind leaves the machine. The cost, control, and privacy objections of Chapter 1 are answered by construction rather than by policy.

\textbf{The schema diagnostic: what fine-tuning actually teaches.} One measurement seals the chapter's argument about where competence lives. With guided decoding \emph{disabled} (n=100, development set), the untrained Qwen3-8B base parses as JSON on 100/100 questions. It produces plans grounded in the question (real entities, real years, real CPV codes, recognisable filter structure) on \textbf{98/100}. But it conforms to the compiler's target contract on only \textbf{3/100}. Opening the raw outputs (deliberately, not trusting the gate) shows the base emitting question-grounded plans in a fluent \emph{self-invented} schema --- \texttt{op:filter\_records}, \texttt{args.filters}, \texttt{id/inputs} --- that the deterministic compiler cannot consume. The SFT and DPO students conform on 98 and 99 respectively. The conclusion is sharp. Fine-tuning does not teach planning ability, because the base already extracts entities and builds filter structure. It teaches \textbf{conformance to the specific executable contract} (3 $\rightarrow$ 98/99). And since guided decoding enforces exactly that contract at evaluation time, the accuracy ladder isolates planning \emph{quality within the contract}, not the ability to emit it. Two recorded caveats bound the claim. The diagnostic stored only failing raws, so trained-arm passes were not individually inspected; the claim rests on the base arm's verified grounded-but-off-schema pattern. And one base abstention expressed through \texttt{return.operation} was under-counted, to which the 3-versus-98 direction is insensitive. The diagnostic itself survived three prior opposite-signed gate flaws, each caught by raw inspection; that audit trail belongs to the evaluation methodology of Chapter 4.

The three layers can now be restated as what the code does. The hard-verification layer pins every projectable property (checks 1--3, with the oracle confined to training time). The abstention layer catches what projection cannot see, gated on provenance so that emptiness is an answer. The learning layer consumes only what the first layer certifies: the oracle-filtered, never oracle-authored, training pools whose consequences the next chapter measures.
